\documentclass[aspectratio=169,11pt]{beamer}

% Vietnamese-friendly Unicode compilation with XeLaTeX or LuaLaTeX
\usepackage{fontspec}
\setsansfont{DejaVu Sans}
\setmonofont{DejaVu Sans Mono}[Scale=0.86]
\usepackage{booktabs}
\usepackage{array}
\usepackage{listings}
\usepackage{xcolor}
\usepackage{xurl}
\usepackage{tikz}
\usetikzlibrary{arrows.meta,positioning,shapes.geometric}

\usetheme{Madrid}
\usecolortheme{seahorse}
\setbeamertemplate{navigation symbols}{}
\setbeamertemplate{footline}[frame number]

\definecolor{sqlblue}{RGB}{28,74,135}
\definecolor{sqlgreen}{RGB}{24,117,74}
\definecolor{sqlgray}{RGB}{245,247,250}
\definecolor{sqlred}{RGB}{166,45,45}

\lstdefinestyle{sqlstyle}{
  language=SQL,
  basicstyle=\ttfamily\scriptsize,
  keywordstyle=\color{sqlblue}\bfseries,
  commentstyle=\color{sqlgreen},
  stringstyle=\color{sqlred},
  showstringspaces=false,
  breaklines=true,
  columns=fullflexible,
  backgroundcolor=\color{sqlgray},
  frame=single,
  rulecolor=\color{gray!35},
  xleftmargin=1mm,
  xrightmargin=1mm,
  aboveskip=4pt,
  belowskip=4pt
}
\lstset{style=sqlstyle}

\title[SQL SELECT]{Giới thiệu câu lệnh \texttt{SELECT} trong SQL}
\subtitle{Truy vấn dữ liệu một bảng với cơ sở dữ liệu mẫu \texttt{classicmodels}}
\author{DBMS Course}
\institute{Smart Logistics and Supply Chain Management Lab}
\date{\today}

\newcommand{\kw}[1]{\textcolor{sqlblue}{\texttt{#1}}}
\newcommand{\minihead}[1]{\vspace{2pt}\textbf{#1}\vspace{2pt}}

\begin{document}

\begin{frame}
  \titlepage
\end{frame}

\begin{frame}{Mục tiêu học tập}
Sau bài này, sinh viên có thể:
\begin{enumerate}
  \item Hiểu vai trò của câu lệnh \kw{SELECT} trong SQL.
  \item Lấy toàn bộ dữ liệu hoặc một số cột cụ thể từ một bảng.
  \item Đặt bí danh cột, tạo cột tính toán trong kết quả truy vấn.
  \item Lọc dữ liệu với \kw{WHERE} và các toán tử điều kiện.
  \item Sắp xếp, giới hạn số dòng, và loại bỏ giá trị trùng lặp.
  \item Viết truy vấn đơn giản trên cơ sở dữ liệu \texttt{classicmodels}.
\end{enumerate}
\end{frame}

\begin{frame}[fragile]{SELECT dùng để làm gì?}
\begin{block}{Ý tưởng chính}
\kw{SELECT} dùng để lấy dữ liệu từ một hoặc nhiều bảng và hiển thị kết quả dưới dạng bảng gồm các dòng và cột.
\end{block}
\begin{columns}[T]
\column{0.48\textwidth}
\minihead{Câu hỏi dữ liệu}
\begin{itemize}
  \item Khách hàng nào ở USA?
  \item Sản phẩm nào có giá trên 150?
  \item Đơn hàng nào đã bị hủy?
\end{itemize}
\column{0.48\textwidth}
\minihead{Trả lời bằng SQL}
\begin{lstlisting}
SELECT customerName, country
FROM customers
WHERE country = 'USA';
\end{lstlisting}
\end{columns}
\end{frame}

\begin{frame}{Cơ sở dữ liệu mẫu: classicmodels}
\begin{itemize}
  \item \texttt{classicmodels} mô phỏng hoạt động của một công ty bán các mô hình xe, tàu, máy bay và sản phẩm sưu tầm.
  \item Trong bài này, ta tập trung vào truy vấn một bảng.
\end{itemize}
\vspace{4pt}
\begin{center}
\begin{tabular}{ll}
\toprule
\textbf{Bảng} & \textbf{Ý nghĩa} \\
\midrule
\texttt{products} & Thông tin sản phẩm \\
\texttt{customers} & Thông tin khách hàng \\
\texttt{orders} & Đơn hàng \\
\texttt{orderdetails} & Chi tiết đơn hàng \\
\texttt{payments} & Thanh toán \\
\texttt{employees} & Nhân viên \\
\texttt{offices} & Văn phòng \\
\texttt{productlines} & Dòng sản phẩm \\
\bottomrule
\end{tabular}
\end{center}
\end{frame}

\begin{frame}[fragile]{Chuẩn bị môi trường thực hành}
\begin{columns}[T]
\column{0.50\textwidth}
\minihead{Chọn cơ sở dữ liệu}
\begin{lstlisting}
USE classicmodels;
\end{lstlisting}
\minihead{Xem danh sách bảng}
\begin{lstlisting}
SHOW TABLES;
\end{lstlisting}
\column{0.50\textwidth}
\minihead{Xem cấu trúc bảng}
\begin{lstlisting}
DESC products;
DESC customers;
DESC orders;
\end{lstlisting}
\begin{alertblock}{Lưu ý}
Trước khi viết truy vấn, nên biết bảng có những cột nào và kiểu dữ liệu của từng cột.
\end{alertblock}
\end{columns}
\end{frame}

\begin{frame}[fragile]{Cú pháp cơ bản của SELECT}
\begin{block}{Mẫu tổng quát}
\begin{lstlisting}
SELECT column1, column2, ...
FROM table_name;
\end{lstlisting}
\end{block}
\begin{center}
\begin{tabular}{ll}
\toprule
\textbf{Thành phần} & \textbf{Ý nghĩa} \\
\midrule
\kw{SELECT} & Chỉ định các cột muốn hiển thị \\
\texttt{column1, column2, ...} & Danh sách cột cần lấy \\
\kw{FROM} & Chỉ định bảng chứa dữ liệu \\
\texttt{table\_name} & Tên bảng cần truy vấn \\
\bottomrule
\end{tabular}
\end{center}
\end{frame}

\begin{frame}[fragile]{Ví dụ đầu tiên}
\begin{columns}[T]
\column{0.52\textwidth}
\minihead{Lấy tên khách hàng}
\begin{lstlisting}
SELECT customerName
FROM customers;
\end{lstlisting}
\minihead{Lấy mã đơn hàng}
\begin{lstlisting}
SELECT orderNumber
FROM orders;
\end{lstlisting}
\column{0.44\textwidth}
\begin{block}{Cách đọc truy vấn}
\begin{itemize}
  \item \kw{SELECT}: lấy cột nào?
  \item \kw{FROM}: từ bảng nào?
  \item Dấu \texttt{;} kết thúc câu lệnh.
\end{itemize}
\end{block}
\end{columns}
\end{frame}

\begin{frame}[fragile]{Lấy toàn bộ cột với SELECT *}
\begin{lstlisting}
SELECT *
FROM productlines;
\end{lstlisting}
\begin{columns}[T]
\column{0.48\textwidth}
\begin{block}{Khi nào hữu ích?}
\begin{itemize}
  \item Xem nhanh dữ liệu trong bảng.
  \item Khám phá bảng khi mới học SQL.
\end{itemize}
\end{block}
\column{0.48\textwidth}
\begin{alertblock}{Không nên lạm dụng}
Với bảng lớn, \kw{SELECT *} có thể lấy quá nhiều cột, khó đọc và tốn tài nguyên. Trong thực tế, nên liệt kê cột cần dùng.
\end{alertblock}
\end{columns}
\end{frame}

\begin{frame}[fragile]{Lấy một số cột cụ thể}
\begin{columns}[T]
\column{0.48\textwidth}
\minihead{Sản phẩm}
\begin{lstlisting}
SELECT productCode, productName, productLine, MSRP
FROM products;
\end{lstlisting}
\minihead{Khách hàng}
\begin{lstlisting}
SELECT customerNumber, customerName, city, country
FROM customers;
\end{lstlisting}
\column{0.48\textwidth}
\begin{block}{Nguyên tắc}
Chỉ lấy các cột cần thiết cho câu hỏi phân tích.
\end{block}
\begin{exampleblock}{Câu hỏi}
Muốn lập danh sách khách hàng theo quốc gia, ta cần cột nào?
\end{exampleblock}
\end{columns}
\end{frame}

\begin{frame}[fragile]{Đặt bí danh cột với AS}
\begin{lstlisting}
SELECT
    productCode AS MaSanPham,
    productName AS TenSanPham,
    productLine AS DongSanPham,
    MSRP AS GiaBanDeXuat
FROM products;
\end{lstlisting}
\begin{itemize}
  \item Bí danh giúp tên cột trong kết quả dễ đọc hơn.
  \item Có thể bỏ từ khóa \kw{AS}, nhưng khi mới học nên dùng \kw{AS} để truy vấn rõ ràng.
\end{itemize}
\end{frame}

\begin{frame}[fragile]{Tạo cột tính toán trong SELECT}
\begin{columns}[T]
\column{0.56\textwidth}
\begin{lstlisting}
SELECT
    productCode,
    productName,
    buyPrice,
    MSRP,
    MSRP - buyPrice AS profitEstimate
FROM products;
\end{lstlisting}
\column{0.40\textwidth}
\begin{block}{Ý nghĩa}
\begin{itemize}
  \item Tính toán ngay trong phần \kw{SELECT}.
  \item Cột tính toán chỉ xuất hiện trong kết quả truy vấn.
  \item Không tự động tạo thêm cột mới trong bảng gốc.
\end{itemize}
\end{block}
\end{columns}
\vspace{4pt}
\begin{lstlisting}
SELECT quantityOrdered * priceEach AS lineTotal
FROM orderdetails;
\end{lstlisting}
\end{frame}

\begin{frame}[fragile]{Lọc dữ liệu với WHERE}
\begin{block}{Mẫu tổng quát}
\begin{lstlisting}
SELECT column1, column2
FROM table_name
WHERE condition;
\end{lstlisting}
\end{block}
\begin{columns}[T]
\column{0.48\textwidth}
\minihead{Sản phẩm Motorcycles}
\begin{lstlisting}
SELECT productCode, productName, productLine
FROM products
WHERE productLine = 'Motorcycles';
\end{lstlisting}
\column{0.48\textwidth}
\minihead{Khách hàng ở USA}
\begin{lstlisting}
SELECT customerNumber, customerName, country
FROM customers
WHERE country = 'USA';
\end{lstlisting}
\end{columns}
\end{frame}

\begin{frame}{Các toán tử so sánh thường dùng}
\begin{center}
\begin{tabular}{lll}
\toprule
\textbf{Toán tử} & \textbf{Ý nghĩa} & \textbf{Ví dụ} \\
\midrule
\texttt{=} & Bằng & \texttt{country = 'USA'} \\
\texttt{<>} & Khác & \texttt{status <> 'Shipped'} \\
\texttt{>} & Lớn hơn & \texttt{MSRP > 150} \\
\texttt{<} & Nhỏ hơn & \texttt{quantityInStock < 1000} \\
\texttt{>=} & Lớn hơn hoặc bằng & \texttt{creditLimit >= 100000} \\
\texttt{<=} & Nhỏ hơn hoặc bằng & \texttt{buyPrice <= 50} \\
\bottomrule
\end{tabular}
\end{center}
\begin{exampleblock}{Ghi nhớ}
Điều kiện trong \kw{WHERE} trả về đúng hoặc sai cho từng dòng dữ liệu.
\end{exampleblock}
\end{frame}

\begin{frame}[fragile]{Ví dụ: lọc theo số}
\begin{columns}[T]
\column{0.50\textwidth}
\minihead{Giá bán đề xuất trên 150}
\begin{lstlisting}
SELECT productCode, productName, MSRP
FROM products
WHERE MSRP > 150;
\end{lstlisting}
\column{0.46\textwidth}
\minihead{Tồn kho dưới 1000}
\begin{lstlisting}
SELECT productCode, productName, quantityInStock
FROM products
WHERE quantityInStock < 1000;
\end{lstlisting}
\end{columns}
\vspace{5pt}
\begin{lstlisting}
SELECT customerNumber, customerName, creditLimit
FROM customers
WHERE creditLimit >= 100000;
\end{lstlisting}
\end{frame}

\begin{frame}[fragile]{Kết hợp điều kiện với AND}
\begin{block}{Ý nghĩa}
\kw{AND} yêu cầu tất cả điều kiện đều đúng.
\end{block}
\begin{lstlisting}
SELECT productCode, productName, productLine, MSRP
FROM products
WHERE productLine = 'Classic Cars'
  AND MSRP > 100;
\end{lstlisting}
\begin{lstlisting}
SELECT customerNumber, customerName, country, creditLimit
FROM customers
WHERE country = 'USA'
  AND creditLimit > 100000;
\end{lstlisting}
\end{frame}

\begin{frame}[fragile]{Kết hợp điều kiện với OR}
\begin{block}{Ý nghĩa}
\kw{OR} yêu cầu ít nhất một điều kiện đúng.
\end{block}
\begin{columns}[T]
\column{0.50\textwidth}
\minihead{Khách hàng ở USA hoặc France}
\begin{lstlisting}
SELECT customerNumber, customerName, country
FROM customers
WHERE country = 'USA'
   OR country = 'France';
\end{lstlisting}
\column{0.46\textwidth}
\minihead{Sản phẩm Motorcycles hoặc Planes}
\begin{lstlisting}
SELECT productCode, productName, productLine
FROM products
WHERE productLine = 'Motorcycles'
   OR productLine = 'Planes';
\end{lstlisting}
\end{columns}
\end{frame}

\begin{frame}[fragile]{Phủ định điều kiện với NOT}
\begin{columns}[T]
\column{0.50\textwidth}
\minihead{Dùng NOT}
\begin{lstlisting}
SELECT customerNumber, customerName, country
FROM customers
WHERE NOT country = 'USA';
\end{lstlisting}
\column{0.46\textwidth}
\minihead{Viết tương đương}
\begin{lstlisting}
SELECT customerNumber, customerName, country
FROM customers
WHERE country <> 'USA';
\end{lstlisting}
\end{columns}
\vspace{5pt}
\begin{lstlisting}
SELECT orderNumber, orderDate, status
FROM orders
WHERE status <> 'Shipped';
\end{lstlisting}
\end{frame}

\begin{frame}[fragile]{Tìm kiếm mẫu chuỗi với LIKE}
\begin{block}{Ý nghĩa}
\kw{LIKE} dùng để so khớp chuỗi theo mẫu.
\end{block}
\begin{center}
\begin{tabular}{ll}
\toprule
\textbf{Ký hiệu} & \textbf{Ý nghĩa} \\
\midrule
\texttt{\%} & Đại diện cho 0 hoặc nhiều ký tự \\
\texttt{\_} & Đại diện cho đúng 1 ký tự \\
\bottomrule
\end{tabular}
\end{center}
\begin{lstlisting}
SELECT customerName, city, country
FROM customers
WHERE customerName LIKE '%Mini%';
\end{lstlisting}
\begin{lstlisting}
SELECT productCode, productName
FROM products
WHERE productName LIKE '196%';
\end{lstlisting}
\end{frame}

\begin{frame}[fragile]{Lọc theo tập giá trị với IN}
\begin{block}{Ý nghĩa}
\kw{IN} giúp viết gọn nhiều điều kiện \kw{OR} trên cùng một cột.
\end{block}
\begin{columns}[T]
\column{0.50\textwidth}
\minihead{Dùng OR}
\begin{lstlisting}
WHERE country = 'USA'
   OR country = 'France'
   OR country = 'Germany'
\end{lstlisting}
\column{0.46\textwidth}
\minihead{Dùng IN}
\begin{lstlisting}
WHERE country IN ('USA', 'France', 'Germany')
\end{lstlisting}
\end{columns}
\vspace{4pt}
\begin{lstlisting}
SELECT customerNumber, customerName, country
FROM customers
WHERE country IN ('USA', 'France', 'Germany');
\end{lstlisting}
\end{frame}

\begin{frame}[fragile]{Lọc theo khoảng với BETWEEN}
\begin{block}{Ý nghĩa}
\kw{BETWEEN a AND b} dùng để kiểm tra giá trị nằm trong khoảng từ \texttt{a} đến \texttt{b}.
\end{block}
\begin{lstlisting}
SELECT productCode, productName, MSRP
FROM products
WHERE MSRP BETWEEN 100 AND 150;
\end{lstlisting}
\begin{lstlisting}
SELECT orderNumber, orderDate, status
FROM orders
WHERE orderDate BETWEEN '2005-01-01' AND '2005-12-31';
\end{lstlisting}
\begin{alertblock}{Lưu ý}
Trong SQL, \kw{BETWEEN} thường bao gồm cả hai đầu mút của khoảng.
\end{alertblock}
\end{frame}

\begin{frame}[fragile]{Kiểm tra giá trị NULL}
\begin{block}{NULL là gì?}
\kw{NULL} biểu diễn giá trị chưa biết, thiếu hoặc không áp dụng được. Không so sánh \kw{NULL} bằng \texttt{=}.
\end{block}
\begin{columns}[T]
\column{0.50\textwidth}
\minihead{Tìm dòng có NULL}
\begin{lstlisting}
SELECT customerNumber, customerName, salesRepEmployeeNumber
FROM customers
WHERE salesRepEmployeeNumber IS NULL;
\end{lstlisting}
\column{0.46\textwidth}
\minihead{Tìm dòng không NULL}
\begin{lstlisting}
SELECT customerNumber, customerName, salesRepEmployeeNumber
FROM customers
WHERE salesRepEmployeeNumber IS NOT NULL;
\end{lstlisting}
\end{columns}
\end{frame}

\begin{frame}[fragile]{Loại bỏ trùng lặp với DISTINCT}
\begin{block}{Ý nghĩa}
\kw{DISTINCT} trả về các giá trị khác nhau trong kết quả truy vấn.
\end{block}
\begin{columns}[T]
\column{0.50\textwidth}
\minihead{Có thể trùng}
\begin{lstlisting}
SELECT country
FROM customers;
\end{lstlisting}
\column{0.46\textwidth}
\minihead{Không trùng}
\begin{lstlisting}
SELECT DISTINCT country
FROM customers;
\end{lstlisting}
\end{columns}
\begin{exampleblock}{Câu hỏi}
Công ty có khách hàng ở những quốc gia nào?
\end{exampleblock}
\end{frame}

\begin{frame}[fragile]{Sắp xếp kết quả với ORDER BY}
\begin{block}{Mẫu tổng quát}
\begin{lstlisting}
SELECT column1, column2
FROM table_name
ORDER BY column1 ASC, column2 DESC;
\end{lstlisting}
\end{block}
\begin{columns}[T]
\column{0.50\textwidth}
\minihead{Tăng dần}
\begin{lstlisting}
SELECT productCode, productName, MSRP
FROM products
ORDER BY MSRP ASC;
\end{lstlisting}
\column{0.46\textwidth}
\minihead{Giảm dần}
\begin{lstlisting}
SELECT productCode, productName, MSRP
FROM products
ORDER BY MSRP DESC;
\end{lstlisting}
\end{columns}
\end{frame}

\begin{frame}[fragile]{Giới hạn số dòng với LIMIT}
\begin{block}{Ý nghĩa}
\kw{LIMIT} giới hạn số dòng trả về. Trong MySQL, thường dùng để lấy một số dòng đầu tiên hoặc phân trang.
\end{block}
\begin{columns}[T]
\column{0.50\textwidth}
\minihead{Lấy 5 sản phẩm đầu}
\begin{lstlisting}
SELECT productCode, productName, MSRP
FROM products
LIMIT 5;
\end{lstlisting}
\column{0.46\textwidth}
\minihead{Top 5 sản phẩm đắt nhất}
\begin{lstlisting}
SELECT productCode, productName, MSRP
FROM products
ORDER BY MSRP DESC
LIMIT 5;
\end{lstlisting}
\end{columns}
\end{frame}

\begin{frame}{Thứ tự viết và thứ tự xử lý logic}
\begin{columns}[T]
\column{0.47\textwidth}
\begin{block}{Thứ tự viết phổ biến}
\begin{enumerate}
  \item \kw{SELECT}
  \item \kw{FROM}
  \item \kw{WHERE}
  \item \kw{ORDER BY}
  \item \kw{LIMIT}
\end{enumerate}
\end{block}
\column{0.49\textwidth}
\begin{block}{Cách hiểu đơn giản}
\begin{enumerate}
  \item Lấy bảng từ \kw{FROM}.
  \item Giữ lại dòng thỏa \kw{WHERE}.
  \item Chọn cột trong \kw{SELECT}.
  \item Sắp xếp bằng \kw{ORDER BY}.
  \item Lấy số dòng bằng \kw{LIMIT}.
\end{enumerate}
\end{block}
\end{columns}
\end{frame}

\begin{frame}[fragile]{Một truy vấn hoàn chỉnh}
\begin{lstlisting}
SELECT
    productCode AS MaSanPham,
    productName AS TenSanPham,
    productLine AS DongSanPham,
    MSRP AS GiaBanDeXuat
FROM products
WHERE productLine IN ('Classic Cars', 'Motorcycles')
  AND MSRP BETWEEN 100 AND 200
ORDER BY MSRP DESC
LIMIT 10;
\end{lstlisting}
\begin{block}{Câu hỏi truy vấn đang trả lời}
Lấy 10 sản phẩm đắt nhất trong hai dòng \texttt{Classic Cars} và \texttt{Motorcycles}, với giá bán đề xuất từ 100 đến 200.
\end{block}
\end{frame}

\begin{frame}{Lỗi thường gặp khi mới học SELECT}
\begin{center}
\begin{tabular}{p{0.36\textwidth}p{0.52\textwidth}}
\toprule
\textbf{Lỗi} & \textbf{Cách sửa} \\
\midrule
Quên dấu phẩy giữa các cột & Viết \texttt{SELECT col1, col2, col3} \\
Quên \kw{FROM} & Luôn xác định bảng dữ liệu nguồn \\
So sánh chuỗi không có dấu nháy & Viết \texttt{country = 'USA'} \\
So sánh \kw{NULL} bằng \texttt{=} & Dùng \kw{IS NULL} hoặc \kw{IS NOT NULL} \\
Dùng \kw{SELECT *} quá nhiều & Chỉ chọn cột cần thiết \\
\bottomrule
\end{tabular}
\end{center}
\end{frame}

\begin{frame}[fragile]{Bài tập nhanh 1}
Viết truy vấn cho các yêu cầu sau:
\begin{enumerate}
  \item Lấy \texttt{customerNumber}, \texttt{customerName}, \texttt{city}, \texttt{country} từ bảng \texttt{customers}.
  \item Lấy các khách hàng ở \texttt{France} hoặc \texttt{Germany}.
  \item Lấy các sản phẩm có \texttt{MSRP} lớn hơn 150.
  \item Lấy 5 đơn hàng mới nhất theo \texttt{orderDate}.
\end{enumerate}
\pause
\begin{block}{Gợi ý cho câu 4}
\begin{lstlisting}
SELECT orderNumber, orderDate, status
FROM orders
ORDER BY orderDate DESC
LIMIT 5;
\end{lstlisting}
\end{block}
\end{frame}

\begin{frame}[fragile]{Bài tập nhanh 2}
Viết truy vấn cho các yêu cầu sau:
\begin{enumerate}
  \item Lấy danh sách quốc gia khác nhau trong bảng \texttt{customers}.
  \item Lấy các sản phẩm có tên chứa chuỗi \texttt{'Ford'}.
  \item Lấy các khách hàng chưa có nhân viên phụ trách.
  \item Tính \texttt{lineTotal = quantityOrdered * priceEach} trong bảng \texttt{orderdetails}.
\end{enumerate}
\pause
\begin{block}{Gợi ý}
\begin{lstlisting}
SELECT DISTINCT country FROM customers;
SELECT * FROM products WHERE productName LIKE '%Ford%';
SELECT * FROM customers WHERE salesRepEmployeeNumber IS NULL;
\end{lstlisting}
\end{block}
\end{frame}

\begin{frame}{Tóm tắt}
\begin{itemize}
  \item \kw{SELECT} là câu lệnh trung tâm để truy vấn dữ liệu.
  \item \kw{FROM} cho biết bảng dữ liệu nguồn.
  \item \kw{WHERE} lọc dòng theo điều kiện.
  \item \kw{AS} đặt bí danh cho cột; biểu thức trong \kw{SELECT} tạo cột tính toán.
  \item \kw{DISTINCT}, \kw{ORDER BY}, \kw{LIMIT} giúp làm sạch, sắp xếp và giới hạn kết quả.
  \item Khi mới học, hãy đọc truy vấn theo câu hỏi dữ liệu: \textit{lấy cột nào, từ bảng nào, lọc điều kiện gì?}
\end{itemize}
\end{frame}

\begin{frame}{Tài liệu tham khảo}
\small
\begin{itemize}
  \item Tutorial DBMS: \textit{Sử dụng câu lệnh SELECT cơ bản với classicmodels}.\\
  {\tiny\url{https://vudmvn.github.io/dbms-course/lectures/MySQL/querying-data/select/select-statement-1/}}
  \item GeeksforGeeks: \textit{SQL SELECT Query}.\\
  {\tiny\url{https://www.geeksforgeeks.org/sql/sql-select-query/}}
\end{itemize}
\end{frame}

\end{document}
